\section{Robust freezing after cofactor descent}
\label{sec:tpc45-residual-freezing}

We now combine the exact band descent with the face argument of
TPC--44.  The proof is included because the diagonal loss
\(\rho_R\) creates a useful quantitative tradeoff between the width of
the physical band and the logarithmic low-prime cutoff.
Write the inherited endpoint budget as
\begin{equation}
 K_B=X^{1/400+\kappa_X},
 \qquad \kappa_X=o(1).
 \label{eq:tpc45-KB-definition}
\end{equation}

Let \(\widetilde S_i=\sum_w\widetilde b_{i,w}\chi_w\) be the grouped
packet in \eqref{eq:tpc45-grouped-coefficients}; none of its labels is
divisible by a prime in \(\cP_R\).  Fix \(z<\sqrt{N_+}\), so the low
cutoff and the band are disjoint.  Split every descended label uniquely
as
\begin{equation}
 w=a b,
 \qquad \Pplus(a)\le z<\Pminus(b),
 \label{eq:tpc45-descended-low-high-split}
\end{equation}
and let \(\widetilde\cA_z\) be the set of low parts \(a\) that occur.
Put
\begin{equation}
 \widetilde\nu_z:=|\widetilde\cA_z|.
 \label{eq:tpc45-descended-projection-count}
\end{equation}
For a remaining high-prime environment \(\xi=(\varepsilon_p)_{p>z,
p\notin\cP_R}\), define
\begin{equation}
 H_{i,a}(\xi)
 :=\sum_{w_{\le z}=a}
   \widetilde b_{i,w}\chi_{w_{>z}}(\xi),
 \qquad
 W_z(\xi):=\sum_i\sum_{a\in\widetilde\cA_z}|H_{i,a}(\xi)|^2.
 \label{eq:tpc45-descended-Wz}
\end{equation}

\begin{lemma}[Robust descended face]
\label{lem:tpc45-robust-descended-face}
For every \(L\ge1\),
\begin{equation}
 \Pp_\xi\!\left(
  \sup_{\sigma_{p\le z}}
  Z_0(-\one_{\cP_R},\sigma,\xi)
  >L\widetilde\nu_z\widetilde\cD
 \right)
 \le L^{-1}.
 \label{eq:tpc45-robust-descended-tail}
\end{equation}
\end{lemma}

\begin{proof}
Distinct descended labels with the same low part have distinct high
parts.  Walsh orthogonality therefore gives
\begin{equation}
 \E_\xi W_z(\xi)=\widetilde\cD.
 \label{eq:tpc45-descended-W-mean}
\end{equation}
For every assignment \(\sigma\) below \(z\),
\begin{equation}
 \widetilde S_i(\sigma,\xi)
 =\sum_{a\in\widetilde\cA_z}\chi_a(\sigma)H_{i,a}(\xi).
 \label{eq:tpc45-descended-coordinate-grouping}
\end{equation}
Cauchy--Schwarz and then summation in \(i\) give
\begin{equation}
 \sup_\sigma Z_0(-\one_{\cP_R},\sigma,\xi)
 \le\widetilde\nu_zW_z(\xi).
 \label{eq:tpc45-descended-face-Cauchy}
\end{equation}
Markov's inequality and \eqref{eq:tpc45-descended-W-mean} prove the
claim.
\end{proof}

Every untouched descended label is an original label of size at most
\begin{equation}
 X^{\lambda_0+o(1)},
 \qquad
 \lambda_0=1+\frac{10049}{52500}
 =\frac{62549}{52500},
 \label{eq:tpc45-full-descended-label-size}
\end{equation}
while touched labels are shorter by
\eqref{eq:tpc45-touched-label-exponent}.  The standard Rankin argument
for squarefree \(z\)-smooth integers gives, when
\(z=(\log X)^A\) and \(A>1\),
\begin{equation}
 \widetilde\nu_z
 \le X^{\lambda_0(1-1/A)+o(1)}.
 \label{eq:tpc45-full-Rankin-projection}
\end{equation}
For completeness, choose
\(s=1-1/A+(\log\log X)^{-1/2}\) and bound the number of possible low
parts by
\begin{equation}
 X^{(\lambda_0+o(1))s}
 \prod_{p\le z}(1+p^{-s}).
 \label{eq:tpc45-Rankin-product}
\end{equation}
The logarithm of the product is
\(O_A(z^{1-s})=o(\log X)\), which proves
\eqref{eq:tpc45-full-Rankin-projection}.

\begin{theorem}[Band/cutoff tradeoff]
\label{thm:tpc45-band-cutoff-tradeoff}
Suppose
\begin{equation}
 R=X^{\theta+o(1)},
 \qquad 0\le\theta<\frac12,
 \qquad z=(\log X)^A,\qquad A>1,
 \label{eq:tpc45-tradeoff-parameters}
\end{equation}
and
\begin{equation}
 \boxed{
 (\lambda_L-\theta)_+
 +\lambda_0\left(1-\frac1A\right)
 <\frac1{400}.}
 \label{eq:tpc45-tradeoff-condition}
\end{equation}
Let
\begin{equation}
 L_X:=\exp\!\left(
  \frac{\log X}{\sqrt{\log\log X}}
 \right).
 \label{eq:tpc45-LX}
\end{equation}
Then, for all sufficiently large \(X\),
\begin{equation}
 \Pp_\xi\!\left(
  \sup_{\sigma_{p\le z}}
  Z_0(-\one_{\cP_R},\sigma,\xi)>K_B\cD_0
 \right)
 \le L_X^{-1}=o(1).
 \label{eq:tpc45-tradeoff-tail}
\end{equation}
Thus a density-one set of the remaining environments tolerates, at
once, the physical assignment on the whole band \(\cP_R\) and every
assignment below \(z\).
\end{theorem}

\begin{proof}
By \cref{cor:tpc45-grouped-diagonal},
\begin{equation}
 \widetilde\cD
 \le X^{(\lambda_L-\theta)_++o(1)}\cD_0.
 \label{eq:tpc45-tradeoff-diagonal}
\end{equation}
Combine this with \eqref{eq:tpc45-full-Rankin-projection}.  Choose
\(\delta>0\) so that the left side of
\eqref{eq:tpc45-tradeoff-condition} is at most
\(1/400-2\delta\).  The strict gap absorbs both the \(X^{o(1)}\)
factors and \(L_X=X^{o(1)}\), giving
\begin{equation}
 L_X\widetilde\nu_z\widetilde\cD
 \le X^{1/400-\delta}\cD_0
 \le K_B\cD_0
 \label{eq:tpc45-tradeoff-budget}
\end{equation}
for sufficiently large \(X\), with the inherited harmless
subpower convention in \(K_B\).  Apply
\cref{lem:tpc45-robust-descended-face} with \(L=L_X\).
\end{proof}

\begin{corollary}[Clean balanced freezing with the legacy cutoff]
\label{cor:tpc45-clean-balanced-freezing}
Take \(R=\ell_+\), \(A=501/500\), and
\(z=(\log X)^{501/500}\).  Then
\begin{equation}
 \lambda_0\left(1-\frac{500}{501}\right)
 =\frac{62549}{26302500},
 \qquad
 \frac1{400}-\frac{62549}{26302500}
 =\frac{12829}{105210000}>0.
 \label{eq:tpc45-clean-balanced-margin}
\end{equation}
Consequently \eqref{eq:tpc45-tradeoff-tail} holds while the entire
band
\begin{equation}
 \sqrt{N_+}<p\le N_-/\ell_+
 \label{eq:tpc45-clean-balanced-band-explicit}
\end{equation}
is fixed to the physical sign \(-1\).
\end{corollary}

The inequality is open, so a small part of the endpoint budget can be
spent to move beyond the clean cofactor scale while keeping the same
low cutoff.

\begin{corollary}[A budgeted extension beyond balance]
\label{cor:tpc45-budgeted-band-extension}
Let
\begin{equation}
 R=\ell_+X^{-1/10000},
 \qquad A=\frac{501}{500}.
 \label{eq:tpc45-budgeted-extension-parameters}
\end{equation}
Then the total power cost and its remaining margin are
\begin{align}
 \frac1{10000}+\frac{62549}{26302500}
 &=\frac{260717}{105210000},
 \label{eq:tpc45-budgeted-extension-cost}\\
 \frac1{400}-\frac{260717}{105210000}
 &=\frac{577}{26302500}>0.
 \label{eq:tpc45-budgeted-extension-margin}
\end{align}
Hence the old low-prime cutoff remains valid and the physically frozen
band extends to exponent
\begin{equation}
 1-\lambda_L+\frac1{10000}
 =\frac{55021}{105000}
 =0.5240095238\ldots.
 \label{eq:tpc45-budgeted-extension-endpoint}
\end{equation}
\end{corollary}

One may trade more of the low cutoff for a wider band.  For example,
\begin{equation}
 R=\ell_+X^{-1/1000},
 \qquad A=\frac{1001}{1000}
 \label{eq:tpc45-wider-tradeoff-parameters}
\end{equation}
has exact margin
\begin{equation}
 \frac1{400}-\frac1{1000}
 -\frac{62549}{52552500}
 =\frac{65119}{210210000}>0,
 \label{eq:tpc45-wider-tradeoff-margin}
\end{equation}
and reaches band exponent
\(110231/210000=0.5249095238\ldots\).
These points illustrate one continuous frontier; neither is asserted
to be optimal for the actual coefficients.

\subsection{The shortened touched component}

Let \(Z_{\rm touch}\) be the equality energy formed only from atoms
touched by \(\cP_{\rm bal}\), after the band signs have been set to
\(-1\) and equal residual labels have been grouped.  This is a
standalone nonnegative Walsh packet.  Its labels have exponent
\(\lambda_1=36299/52500\) and its grouped diagonal is at most
\(X^{o(1)}\cD_{\rm touch}\), where \(\cD_{\rm touch}\) is its original
atomic diagonal.

\begin{corollary}[Improved cutoff for the touched packet]
\label{cor:tpc45-touched-improved-cutoff}
For
\begin{equation}
 z_1=(\log X)^{1003/1000},
 \qquad
 A_1^*
 :=\left(1-\frac1{400\lambda_1}\right)^{-1}
 =\frac{145196}{144671},
 \label{eq:tpc45-touched-A-star}
\end{equation}
one has \(1<1003/1000<A_1^*\) and
\begin{equation}
 \frac1{400}
 -\lambda_1\left(1-\frac{1000}{1003}\right)
 =\frac{30329}{70210000}>0.
 \label{eq:tpc45-touched-cutoff-gap}
\end{equation}
Consequently a density-one set of the residual high-prime environments
controls every assignment below \(z_1\) for \(Z_{\rm touch}\) within
the \(K_B\cD_{\rm touch}\) scale.
\end{corollary}

This last corollary is intentionally componentwise.  The untouched
column still has labels of exponent \(\lambda_0\), and the norm of the
sum of touched and untouched columns contains cross terms.  Therefore
\((\log X)^{1003/1000}\) is not a full-packet cutoff.

Finally, the failure probability \(L_X^{-1}\) is summable on dyadic
scales.  Prescribe one inherited packet and one admissible pair
\((R_X,z_X)\) at each \(X=2^k\), and couple the remaining signs by a
single global independent prime-sign field.  The first
Borel--Cantelli lemma makes the corresponding tradeoff assertion hold
at all sufficiently large dyadic scales almost surely.  The same is
true for any fixed finite family of such prescribed choices after a
union bound.  No independence between scales is required; no
uniformity over all packet or parameter sequences is claimed.
